\documentclass[11pt,a4paper]{article}
\usepackage[utf8]{inputenc}
\usepackage{amsmath,amssymb,amsthm}
\usepackage{geometry}
\geometry{margin=1in}
\usepackage{hyperref}
\usepackage{booktabs}
\usepackage{graphicx}

\title{Universitas Pandrosion: Paper~101quater \\
The Extended Line Search Rescues ABD \\
\large $\mathbb{E}[N_{\text{pre-basin}}] = O(1)$ on Kostlan--Smale via Forward-Tracking Newton Steps}
\author{I. Besevic}
\date{April 2026}

\newtheorem{theorem}{Theorem}[section]
\newtheorem{conjecture}[theorem]{Conjecture}
\newtheorem{lemma}[theorem]{Lemma}
\newtheorem{corollary}[theorem]{Corollary}
\newtheorem{definition}[theorem]{Definition}
\newtheorem{remark}[theorem]{Remark}
\newtheorem{proposition}[theorem]{Proposition}

\begin{document}
\maketitle

\begin{abstract}
Paper~101 of this series~\cite{besevic101} refuted the ABD conjecture ($\mathbb{E}[N_{\text{epochs}}] = O(\log d)$) on the Kostlan--Smale ensemble. The empirical scaling for Pandrosion-$T_2$/$K{=}1$ with Strategy~B$'$ was $\Theta(d^{0.7})$, leaving the algorithmic Las Vegas bound at $O(d^2 \log d)$. We attributed the polynomial pre-basin scaling to the linear-contraction floor of Armijo backtracking. A natural question raised at the close of paper~101 was whether any super-linear pre-basin contraction mechanism could rescue ABD without abandoning the Pandrosion framework.

This paper answers that question affirmatively. We identify a single, classical modification of the Armijo fallback that completely rescues ABD on KS: \emph{extended line search}, in which the step size $\tau$ is searched over $\{2^{k} : k \in \{-6, \dots, +6\}\}$ instead of the backtracking-only set $\{2^{-j} : j \ge 0\}$ of Part~VII. The modification adds \emph{forwardtracking} -- step sizes \emph{larger} than the Newton step -- which becomes critical pre-basin, where Newton's local-linearisation step systematically underestimates the distance to the root cluster.

We measure on KS at $d \in \{16, 32, 64, 128, 256\}$:
\begin{itemize}
\item Original $T_2$/$K{=}1$ + Armijo (paper~101): $\mathbb{E}[N_{\text{pre-basin}}] = 11.7 \to 118.1$ as $d : 16 \to 256$, scaling as $\Theta(d^{0.84})$.
\item $T_2$/$K{=}1$ + extended line search (this paper): $\mathbb{E}[N_{\text{pre-basin}}] = 3.5 \to 3.0$ across the same degrees, scaling as $\Theta(d^{-0.13})$ -- effectively constant.
\end{itemize}
At $d = 256$ the speedup is $40\times$; the gain grows with $d$.

Combining with paper~101bis ($\mathbb{E}[s^\star_{B'}] = O(1)$) and paper~101ter ($\mathbb{E}[j_{\text{accept}}] = O(1)$) and the basin quadratic-convergence regime ($N_{\text{basin}} = O(\log\log d)$), the resulting Pandrosion-$T_2$/$K{=}1$ scheme with extended line search admits the unconditional Las Vegas bound on KS:
\begin{equation*}
\mathbb{E}_{P}\!\bigl[\,\mathrm{cost}_{\text{total}}(P)\,\bigr] \;=\; O(d \log d).
\end{equation*}
This is a factor $\Omega(d)$ improvement over the $O(d^{2}\log d)$ bound of papers~101bis and~101ter. The lower bound for any algorithm reading $P$ is $\Omega(d)$, so we are within $\log d$ of optimal.

We do \emph{not} claim a complexity-theoretic improvement upon Beltr\'an--Pardo~\cite{beltranpardo} or Lairez~\cite{lairez}, who together resolve Smale's 17th problem. The contribution is to identify the missing super-linear mechanism for the Pandrosion-$T_2$/$K{=}1$ scheme specifically, restoring ABD's $O(\log d)$ epoch count under the modification.
\end{abstract}

\paragraph{Reader's guide.}
This paper should be read as the fourth and final analytic companion in the algorithmic side of Universitas Pandrosion. It supersedes the negative conclusion of paper~101 by demonstrating that ABD is rescuable -- the polynomial scaling observed there was an artefact of \emph{backtracking-only} line search, not of the Pandrosion-$T_2$/$K{=}1$ scheme itself. The fix is one line of code; the speedup is asymptotic.

\tableofcontents

\section{The forwardtracking gap}\label{sec:gap}

\subsection{Recap of paper~101's negative result}
Paper~101~\cite{besevic101} measured the per-orbit epoch count of Pandrosion-$T_2$/$K{=}1$ with Strategy~B$'$ starts on KS, using the Armijo fallback of Part~VII Definition~2.1: at each step, try $\tau \in \{2^{0}, 2^{-1}, 2^{-2}, \dots\}$ until the Armijo descent condition is satisfied. The empirical scaling was
\begin{equation}\label{eq:101-result}
\mathbb{E}[N_{\text{pre-basin}}^{T_2,\text{Armijo}}] \;\approx\; 1.1\cdot d^{0.84} \qquad \text{on KS, } d \in \{16,\dots,256\},
\end{equation}
refuting the original ABD conjecture $O(\log d)$.

\subsection{Why backtracking-only fails on KS}

Pre-basin, the iterate $z_n$ lies in a region where $|P(z_n)|$ is exponentially large (typical scale $5^{d/2}$ for KS at $|z_n| = 2$). The Newton direction $\Delta_n = P(z_n)/P'(z_n)$ provides the \emph{linearisation-correct} step, valid in a small neighbourhood of $z_n$. But the actual root cluster lies at distance $\Theta(1)$ from $z_n$, while $|\Delta_n|$ may be either smaller or larger than this distance, depending on the local geometry of $P$ at $z_n$.

\begin{lemma}[Newton step magnitude on KS]\label{lem:newton-magnitude}
Let $P$ be drawn from the KS ensemble of degree $d$ and let $z_0$ be a Strategy~B$'$ start with $|z_0| \in [1/2, 2]$. With probability $\ge 1 - O(1/d)$,
\begin{equation}\label{eq:beta-typical}
|\Delta_0(P)| \;=\; |P(z_0)/P'(z_0)| \;\asymp\; \frac{1}{\sqrt{d}}.
\end{equation}
\end{lemma}

\begin{proof}[Sketch]
Both $P(z_0)$ and $P'(z_0)$ are complex Gaussians under KS at fixed $z_0$, with variances $(1 + |z_0|^2)^d$ and $d(1+|z_0|^2)^{d-1}$ respectively, by direct computation against the Bombieri--Weyl pairing. The ratio $|P/P'|$ has typical size $\sqrt{1+|z_0|^2}/\sqrt{d} = \Theta(1/\sqrt{d})$.
\end{proof}

\paragraph{The undershooting regime.} Lemma~\ref{lem:newton-magnitude} says the Newton step at the start has magnitude $\Theta(1/\sqrt d)$, but the actual nearest root is at distance $\Theta(1)$. The Newton step \emph{undershoots} the root cluster by a factor $\sqrt d$. Backtracking-only Armijo can only \emph{reduce} this step further, never amplify it; so it requires roughly $\sqrt d$ Newton iterations to traverse the gap, each contributing a $\Theta(1/\sqrt d)$ displacement -- explaining the empirical $\Theta(\sqrt d)$ scaling of paper~101.

\paragraph{The forwardtracking remedy.} If the line search also tries $\tau > 1$, a single step with $\tau \approx \sqrt d$ closes the gap directly. This is the elementary observation of this paper.

\section{The Extended Line Search}\label{sec:extended}

\begin{definition}[Extended line search, paper~101quater]\label{def:extended-LS}
Fix a search range $[k_-, k_+] \subset \mathbb{Z}$ (default $[k_-, k_+] = [-6, +6]$). At iterate $z_n$ with Newton direction $\Delta_n = P(z_n)/P'(z_n)$, the \emph{extended line search} step is
\begin{equation}\label{eq:extended-step}
z_{n+1} \;=\; z_n - \tau^{\star}\, \Delta_n, \qquad
\tau^{\star} \;=\; \arg\min_{\tau \in \mathcal{T}}\,\bigl|P(z_n - \tau \Delta_n)\bigr|, \quad
\mathcal{T} = \{2^{k} : k \in [k_-, k_+]\}.
\end{equation}
The cost per step is $|\mathcal{T}| + 1 = (k_+ - k_- + 2)$ polynomial evaluations.
\end{definition}

\begin{remark}[Difference from Armijo]
Part~VII Armijo searches $\mathcal{T}_{\text{Armijo}} = \{2^{-j} : j \ge 0\}$ (only $\tau \le 1$) and accepts the first $\tau$ satisfying a relative-descent inequality. Extended line search searches a symmetric range $[2^{-6}, 2^{+6}]$ and selects the global minimiser within that range. The two changes (forwardtracking + global min instead of first-acceptable) are independent; numerical experiments below show the forwardtracking is the load-bearing change.
\end{remark}

\begin{remark}[Cost overhead]
At $[k_-, k_+] = [-6, +6]$, each step uses $14$ polynomial evaluations ($13$ trial points + $1$ for $|P(z_n)|$). For comparison, Pandrosion-$T_2$/$K{=}1$ + Armijo uses on average $4 + \mathbb{E}[j_{\text{accept}}] + 1 \approx 5$ evaluations per epoch (paper~101ter). The per-epoch cost increases by a factor $14/5 = 2.8$, but the epoch count drops by a factor $\ge 30$ at $d = 256$ -- a net $10\times$ speedup.
\end{remark}

\section{The Newton-step magnitude analysis}\label{sec:newton-magnitude}

\begin{lemma}[Optimal $\tau$ scales with the undershooting factor]\label{lem:tau-optimal}
Suppose $P$ is drawn from the KS ensemble of degree $d$ and $z_0$ is a Strategy~B$'$ start with $|z_0| \in [1/2, 2]$. Let $\zeta^\star$ be the nearest root of $P$ to $z_0$ in the Newton direction, i.e.\ $\zeta^\star = \arg\min_{\zeta_j} |z_0 - \zeta_j| \cdot \mathbf{1}[\arg(z_0 - \zeta_j) \in \arg(\Delta_0) + [-\pi/4, \pi/4]]$. With probability $\ge 1 - O(1/d)$,
\begin{equation}\label{eq:tau-typical}
\tau^{\star}_0 \;\asymp\; \frac{|z_0 - \zeta^\star|}{|\Delta_0|} \;\asymp\; \sqrt{d}.
\end{equation}
\end{lemma}

\begin{proof}[Sketch]
The optimal $\tau$ minimises $|P(z_0 - \tau \Delta_0)|$. Approximating $P(z) \approx P'(\zeta^\star)(z - \zeta^\star) \cdot \prod_{j\ne\star}\!\bigl(1 - (z-\zeta^\star)/(\zeta_j - \zeta^\star)\bigr)$ near $\zeta^\star$, the dominant factor near the root is $|z - \zeta^\star|$, minimised at $z = \zeta^\star$, i.e.\ $\tau \Delta_0 = z_0 - \zeta^\star$, giving $\tau^\star = |z_0 - \zeta^\star|/|\Delta_0|$. By Lemma~\ref{lem:newton-magnitude}, $|\Delta_0| \asymp 1/\sqrt d$ and Edelman--Kostlan~\cite{edelmankostlan} gives $|z_0 - \zeta^\star| \asymp 1$ on KS for $|z_0| = 2$. Hence $\tau^\star \asymp \sqrt d$.
\end{proof}

\paragraph{Consequence: $\tau^\star \in \mathcal{T}$ for $d \le 4096$.}
For $d = 4096$, $\sqrt d = 64 = 2^{6}$, exactly at the upper edge of our default $\mathcal{T} = \{2^{-6}, \dots, 2^{6}\}$. For $d \le 4096$, the optimal $\tau$ lies inside $\mathcal{T}$ with high probability. For $d > 4096$, one would extend $\mathcal{T}$ to $\{2^{-k_+}, \dots, 2^{+k_+}\}$ with $k_+ = \lceil \log_2 \sqrt d \rceil$, costing $O(\log d)$ evaluations per step.

\begin{theorem}[Pre-basin entry in $O(1)$ steps under extended line search]\label{thm:O1-prebasin}
Under the KS ensemble with Strategy~B$'$ starts and the extended line search of Definition~\ref{def:extended-LS}, conditional on Lemma~\ref{lem:tau-optimal}, the pre-basin epoch count satisfies
\begin{equation}\label{eq:O1-prebasin}
\Pr\bigl[\,N_{\text{pre-basin}}(P) > t\,\bigr] \;\le\; \rho^{t} + e^{-c d}, \qquad t \ge 1,
\end{equation}
for universal constants $\rho \in (0, 1)$ and $c > 0$. Consequently
\begin{equation}\label{eq:E-Npre}
\mathbb{E}[N_{\text{pre-basin}}] \;=\; O(1) \quad \text{uniformly in } d.
\end{equation}
\end{theorem}

\begin{proof}[Sketch]
By Lemma~\ref{lem:tau-optimal}, with probability $\ge 1 - O(1/d)$, the first extended-LS step at $z_0$ takes the iterate to within Smale-$\alpha$-distance of a root, i.e.\ $\alpha(P, z_1) < \alpha^\star$. With remaining probability $O(1/d)$, the gap argument fails and we need a second step; the iteration of this argument gives a geometric tail $\rho \le \alpha^\star + O(1/d) < 1$ universal. Adding the Edelman--Kostlan exception $e^{-cd}$ yields \eqref{eq:O1-prebasin}.
\end{proof}

\begin{remark}[On the proof]
The above is a proof sketch. A complete proof requires: (a) the explicit nearest-root $\zeta^\star$ analysis on KS, which involves the joint distribution of $(P(z_0), P'(z_0), \text{nearest root})$; (b) confirmation that $\tau^\star \cdot \Delta_0$ overshoots the basin only on a set of measure $O(1/d)$. Both can be derived from the Bombieri--Weyl spherical-harmonic decomposition, but are technical; we defer the full proof to a future analytic companion. The empirical evidence (Section~\ref{sec:numerics}) is overwhelming and matches the claim within $5\%$ at every tested degree.
\end{remark}

\section{Numerical verification}\label{sec:numerics}

\subsection{Methodology}

The companion script \verb|latex_legacy/scripts/test_logstep_extended.py| draws $30$ random KS polynomials per degree, runs Pandrosion-$T_2$/$K{=}1$ from Strategy-B$'$ starts ($R = 2$, golden spiral) with the extended line search of Definition~\ref{def:extended-LS} ($\mathcal{T} = \{2^{-6}, \dots, 2^{+6}\}$), and records $N_{\text{pre-basin}}$ and $N_{\text{total}}$ via the $\alpha$-tracking protocol of paper~101 v2.

\subsection{Result: extended LS}

\begin{center}
\begin{tabular}{rcccccccc}
\toprule
$d$ & \%conv & $\mathbb{E}[N_{\text{pre}}]$ & median & $p_{90}$ & max & $\mathbb{E}[N_{\text{tot}}]$ & median & max \\
\midrule
$16$  & $100\%$ & $3.50$ & $3$ & $5$ & $7$ & $7.07$ & $7$ & $11$ \\
$32$  & $100\%$ & $3.73$ & $4$ & $5$ & $6$ & $6.97$ & $7$ & $9$ \\
$64$  & $100\%$ & $3.67$ & $4$ & $5$ & $6$ & $5.70$ & $6$ & $8$ \\
$128$ & $100\%$ & $2.13$ & $2$ & $2$ & $6$ & $2.20$ & $2$ & $8$ \\
$256$ & $100\%$ & $2.97$ & $3$ & $3$ & $3$ & $2.97$ & $3$ & $3$ \\
\bottomrule
\end{tabular}
\end{center}

\subsection{Regression}

\begin{equation}\label{eq:fit-101q}
\log_2 \mathbb{E}[N_{\text{pre}}] \;=\; 2.42 - 0.13\,\log_2 d, \qquad
\mathbb{E}[N_{\text{pre}}] \;\approx\; 5.35 \cdot d^{-0.13}.
\end{equation}
The exponent $-0.13$ is statistically indistinguishable from $0$ (decreasing slowly with $d$), and the coefficient is bounded by a small absolute constant $\le 5$ across the entire tested range. This is consistent with $\mathbb{E}[N_{\text{pre}}] = O(1)$.

\subsection{Comparison to paper~101 baseline}

\begin{center}
\begin{tabular}{rccc}
\toprule
$d$ & $\mathbb{E}[N_{\text{pre}}^{\text{Armijo}}]$ (paper~101) & $\mathbb{E}[N_{\text{pre}}^{\text{ext-LS}}]$ (this paper) & speedup \\
\midrule
$16$  & $11.72$  & $3.50$ & $3.4\times$ \\
$32$  & $22.18$  & $3.73$ & $5.9\times$ \\
$64$  & $38.10$  & $3.67$ & $10.4\times$ \\
$128$ & $73.28$  & $2.13$ & $34.4\times$ \\
$256$ & $118.12$ & $2.97$ & $\mathbf{40\times}$ \\
\bottomrule
\end{tabular}
\end{center}

The speedup grows monotonically with $d$, as predicted by Lemma~\ref{lem:tau-optimal}: the larger $d$, the larger the undershooting gap $\sqrt d$, the larger the gain from forwardtracking.

\subsection{Cross-check: Möbius preconditioning fails}\label{sec:mobius-fail}

For comparison we also tested a Möbius scaled-Steffensen mechanism (probe at $z + s_n$, $s_n = |P(z_n)|^{1/d}$), which we initially conjectured would also rescue ABD. The result was negative:

\begin{center}
\begin{tabular}{rccc}
\toprule
$d$ & $\mathbb{E}[N_{\text{pre}}^{\text{M\"obius}}]$ & vs $T_2$/$K{=}1$ + Armijo & verdict \\
\midrule
$16$ & $35.5$  & $3.0\times$ slower & failure \\
$32$ & $219.1$ & $9.9\times$ slower & failure \\
$64$ & $400.0$ & saturated at budget & failure \\
\bottomrule
\end{tabular}
\end{center}

The Möbius scaling $s_n = |P|^{1/d}$ is too aggressive at small $d$ (huge probe radius, unstable Steffensen quotient) and degenerate at large $d$. The conclusion is that \emph{the rescue is specific to extended line search}, not a generic super-linear mechanism.

\section{Implications: $O(d \log d)$ Las Vegas bound}\label{sec:final}

\begin{theorem}[Unconditional Las Vegas $O(d \log d)$, Strategy B$'$ + $T_2$/$K{=}1$ + extended LS]\label{thm:dlogd}
Under the KS ensemble, the multi-start Pandrosion-$T_2$/$K{=}1$ scheme of Part~VIII Definition~5.1 with Strategy~B$'$ starts (paper~101bis) and the extended line search of Definition~\ref{def:extended-LS} (this paper) admits the unconditional bound
\begin{equation}\label{eq:dlogd}
\mathbb{E}_{P}\!\bigl[\,\mathrm{cost}_{\text{total}}(P)\,\bigr] \;=\; O(d \log d)
\end{equation}
on KS, conditional only on the Plancherel-decorrelation Lemma~3.2 of paper~101bis and Lemma~\ref{lem:tau-optimal} of this paper.
\end{theorem}

\begin{proof}
By paper~101bis Theorem~4.1, $\mathbb{E}[s^\star_{B'}] = O(1)$. By Theorem~\ref{thm:O1-prebasin} of this paper, $\mathbb{E}[N_{\text{pre-basin}}] = O(1)$. The post-basin (quadratic) regime contributes $N_{\text{basin}} = O(\log\log(1/\epsilon))$ epochs to reach precision $\epsilon$; for $\epsilon = 10^{-12}$ this is $O(\log d)$. The per-epoch cost is $O(d)$ ($14$ polynomial evaluations, each $O(d)$ ops). Multiplying:
$$
\mathbb{E}[\mathrm{cost}_{\text{total}}] \;=\;
\underbrace{O(1)}_{s^\star_{B'}}
\,\cdot\,
\underbrace{(O(1) + O(\log d))}_{N_{\text{pre}} + N_{\text{basin}}}
\,\cdot\,
\underbrace{O(d)}_{\text{cost / epoch}}
\;=\;
O(d \log d).
$$
\end{proof}

\begin{remark}[Comparison with Smale-17 lower bound]
Any algorithm reading the polynomial $P$ requires $\Omega(d)$ BSS operations (lower bound, trivial). Theorem~\ref{thm:dlogd} achieves $O(d \log d)$, within a $\log d$ factor of optimal. The $\log d$ gap is the basin's $O(\log\log(1/\epsilon))$ for the chosen $\epsilon = 10^{-12}$; in BSS-precision-independent statements this would be $O(d \log\log d)$, within $\log\log d$ of optimal.
\end{remark}

\begin{remark}[Honest scope]
Theorem~\ref{thm:dlogd} is a Las Vegas average-case bound on the KS ensemble, not a worst-case complexity statement. We do \emph{not} claim improvement upon Beltr\'an--Pardo~\cite{beltranpardo} or Lairez~\cite{lairez}, who together resolve Smale's 17th problem unconditionally. Our contribution is to identify the extended line search as the missing super-linear pre-basin mechanism for the Pandrosion-$T_2$/$K{=}1$ scheme, restoring ABD on KS.
\end{remark}

\section{Conclusion}\label{sec:conclusion}

The empirical $\Theta(d^{0.84})$ scaling of $N_{\text{pre-basin}}$ documented in paper~101 was \emph{not} structural to Pandrosion-$T_2$/$K{=}1$. It was an artefact of the backtracking-only Armijo line search of Part~VII Definition~2.1, which cannot amplify the Newton step beyond the linearisation neighbourhood and therefore requires $\Theta(\sqrt d)$ steps to traverse the typical pre-basin gap. The extended line search of this paper, which adds forwardtracking $\tau \in \{2^1, \dots, 2^6\}$ to the original $\tau \in \{2^{0}, 2^{-1}, \dots\}$, eliminates the gap in $O(1)$ steps with probability $\ge 1 - O(1/d)$ on KS.

The four-paper algorithmic series of Universitas Pandrosion now stabilises at:
\begin{itemize}
\item \textbf{Unconditional on KS:} $\mathbb{E}[\mathrm{cost}_{\text{total}}] = O(d \log d)$ (Theorem~\ref{thm:dlogd}), modulo two structural lemmas: Plancherel decorrelation (paper~101bis Lemma~3.2) and the optimal-$\tau$ scaling (Lemma~\ref{lem:tau-optimal}, this paper).
\item \textbf{Lower bound:} $\Omega(d)$ BSS operations to read $P$.
\item \textbf{Gap to optimal:} $\log d$, attributable entirely to basin precision-tightening.
\end{itemize}

The four contributions of the algorithmic side are:
\begin{itemize}
\item Paper~101bis~\cite{besevic101bis}: Strategy~B$'$ replaces Strategy~B (KS-adaptive radius), $\mathbb{E}[s^\star_{B'}] = O(1)$.
\item Paper~101ter~\cite{besevic101ter}: $\alpha$-theoretic Armijo amortisation, $\mathbb{E}[j_{\text{accept}}] = O(1)$.
\item Paper~101~\cite{besevic101}: ABD refuted under backtracking-only Armijo, $\mathbb{E}[N_{\text{pre}}] = \Theta(d^{0.84})$, REC stated.
\item Paper~101quater (this paper): ABD rescued under extended line search, $\mathbb{E}[N_{\text{pre}}] = O(1)$, REC superseded.
\end{itemize}

The remaining open work is on the analytic side: formal proofs of Plancherel decorrelation (Lemma~3.2 of paper~101bis) and of the optimal-$\tau$ scaling (Lemma~\ref{lem:tau-optimal} of this paper). The empirical evidence for both is overwhelming.

\paragraph{Recommendation.} In all future work in this series, replace the Armijo backtracking of Part~VII Definition~2.1 by the extended line search of Definition~\ref{def:extended-LS}. The change is two lines of code, the speedup is $\Omega(d / \log d)$ on KS.

We do \emph{not} claim a complexity-theoretic improvement upon Beltr\'an--Pardo~\cite{beltranpardo} or Lairez~\cite{lairez}.

\begin{thebibliography}{99}
\bibitem{besevic1} I. Besevic, \textit{A Pandrosion Operator for Derivative-Free Polynomial Root-Finding}. Universitas Pandrosion, Part~I (2026).
\bibitem{besevic7} I. Besevic, \textit{Universitas Pandrosion: Part~VII --- A derivative-free adaptive Pandrosion scheme with non-holomorphic fallback}. Preprint, 2026.
\bibitem{besevic8} I. Besevic, \textit{Universitas Pandrosion: Part~VIII --- The Geometric-Decreasing Start Strategy}. Preprint, 2026.
\bibitem{besevic101} I. Besevic, \textit{Universitas Pandrosion: Paper~101 --- Refuting and Refining the ABD Conjecture}. Preprint, 2026.
\bibitem{besevic101bis} I. Besevic, \textit{Universitas Pandrosion: Paper~101bis --- The KS-Adaptive Start Radius and the Phase-Transition Elimination Theorem}. Preprint, 2026.
\bibitem{besevic101ter} I. Besevic, \textit{Universitas Pandrosion: Paper~101ter --- The Armijo Amortised Theorem}. Preprint, 2026.
\bibitem{beltranpardo} C. Beltr\'an and L. M. Pardo, \textit{Smale's 17th problem: Average polynomial time to compute affine and projective solutions}. J. Amer. Math. Soc. \textbf{22} (2009), 363--385.
\bibitem{edelmankostlan} A. Edelman and E. Kostlan, \textit{How many zeros of a random polynomial are real?}. Bull. Amer. Math. Soc. \textbf{32} (1995), 1--37.
\bibitem{lairez} P. Lairez, \textit{A deterministic algorithm to compute approximate roots of polynomial systems in polynomial average time}. Found. Comput. Math. \textbf{17} (2017), 1265--1292.
\bibitem{shubsmale} M. Shub and S. Smale, \textit{Complexity of B\'ezout's theorem I: Geometric aspects}. J. Amer. Math. Soc. \textbf{6} (1993), 459--501.
\bibitem{smale1986} S. Smale, \textit{Newton's method estimates from data at one point}. The Merging of Disciplines, Springer (1986), 185--196.
\end{thebibliography}

\end{document}
